\section*{Chapitre \ref{ch:b3:rate-theories} --- Théories des vitesses de réaction}

\begin{solution}{exo:b3:rate-theories:1}
$\bar v_{\mathrm{rel}} = \sqrt{8 \times 1.381\times10^{-23} \times 298/(\pi \times 18.46 \times 1.661\times10^{-27})} = \qty{585}{m/s}$.
$\sigma\bar v_{\mathrm{rel}}N_A = 0.43\times10^{-18} \times 585 \times 6.022\times10^{23} = \qty{1.51e8}{m^3.mol^{-1}.s^{-1}} =
\qty{1.5e11}{L.mol^{-1}.s^{-1}}$. Valeur mesurée~: $2.07\times10^{-12} \times 6.022\times10^{23}\times10^{-3} =
\qty{1.25e9}{L.mol^{-1}.s^{-1}}$, donc $P = 0.008$~: moins d'une
collision sur cent a la bonne orientation (l'atome O de NO doit
rencontrer un O terminal de \ce{O3}).
\end{solution}

\begin{solution}{exo:b3:rate-theories:2}
$\Delta^{\ddagger}H^\circ = 75.0 - 2.48 = \qty{72.5}{kJ/mol}$ en solution~; $75.0 - 4.96 = \qty{70.0}{kJ/mol}$
pour une réaction bimoléculaire en phase gazeuse.
\end{solution}

\begin{solution}{exo:b3:rate-theories:3}
Diels--Alder~: fortement négative (deux molécules forment un seul
complexe cyclique). Perte de \ce{N2}~: positive (une liaison se relâche,
de la liberté est gagnée). Ions créés dans un solvant polaire~:
négative, les charges qui apparaissent ordonnent le solvant autour
d'elles.
\end{solution}

\begin{solution}{exo:b3:rate-theories:4}
Eau~: $8 \times 8.314 \times 298.15/(3 \times \num{0.890e-3}) = \qty{7.4e9}{L.mol^{-1}.s^{-1}}$~;
hexane
\qty{2.2e10}{L.mol^{-1}.s^{-1}}. L'hexane, trois fois moins visqueux~:
les molécules y diffusent trois fois plus vite.
\end{solution}

\begin{solution}{exo:b3:rate-theories:5}
$\ln(k_2T_1/k_1T_2) = \ln(10 \times 298.15/318.15) = 2.238$~; $\Delta^{\ddagger}H^\circ = 8.3145 \times 2.238/(1/298.15 -
1/318.15) = \qty{88.2}{kJ/mol}$. $\Delta^{\ddagger}S^\circ = R[\ln(kh/k_BT) + \Delta^{\ddagger}H^\circ/RT] = 8.3145 \times
(-36.365 + 35.596) = \qty{-6.4}{J.K^{-1}.mol^{-1}}$ (à \qty{298.15}{K}). $E_a = R\ln10/(1/298.15 - 1/318.15) =
\qty{90.8}{kJ/mol} \approx \Delta^{\ddagger}H^\circ + RT$ à la température
moyenne.
\end{solution}

\begin{solution}{exo:b3:rate-theories:6}
$\mu_{\mathrm{OH}} = 16/17$, $\mu_{\mathrm{OD}} = 32/18$, rapport des nombres d'onde $\sqrt{1.889} = 1.374$~:
$\tilde\nu_{\mathrm{OD}} = \qty{2661}{cm^{-1}}$. $k_{\mathrm H}/k_{\mathrm D} = \exp(1.4388 \times 996/(2 \times 298)) = \eu^{2.40} = 11$.
\end{solution}

\begin{solution}{exo:b3:rate-theories:7}
Valeur maximale $\exp(1.4388 \times 808/(2 \times 298)) = \eu^{1.95} = 7.0$.
Une valeur de 40 ne peut pas venir des \omterm{def:b3:quantum-model-systems:oscillator}{énergies de point zéro}~:
l'hydrogène traverse la barrière par \omterm{def:b3:quantum-model-systems:tunnelling}{effet tunnel}, ce que le deutérium,
deux fois plus lourd, fait bien moins.
\end{solution}

\begin{solution}{exo:b3:rate-theories:8}
$\ce{S2O8^{2-}}$ et \ce{I-}~: $z_{\mathrm A}z_{\mathrm B} = +2$, $\log(k/k_0) = 2 \times 0.51 \times 2 \times 0.10 = 0.20$~: $k$
est multiplié par 1.6. Ester et \ce{OH-}~: un partenaire neutre, pas
d'effet.
$\ce{[Co(NH3)5Br]^{2+}}$ et \ce{OH-}~: $z_{\mathrm A}z_{\mathrm B} = -2$,
$k$ est multiplié par $10^{-0.20} = 0.63$.
\end{solution}

\begin{solution}{exo:b3:rate-theories:9}
Précoce, dans la vallée d'entrée (un état de transition proche des
réactifs, d'après le postulat de Hammond). La réaction inverse,
endothermique, présente une barrière tardive sur son propre chemin~:
c'est l'énergie de vibration de la liaison H--F qui la favorise le
mieux.
\end{solution}

\begin{solution}{exo:b3:rate-theories:10}
$\Delta^{\ddagger}H^\circ = -Rb = \qty{61.9}{kJ/mol}$, d'écart-type $R \times 31 = \qty{0.26}{kJ/mol}$ et
d'intervalle à 95\,\% $\pm\qty{0.8}{kJ/mol}$.
\[
  \Delta^{\ddagger}S^\circ = R(16.528 - 23.760) = \qty{-60.1}{J.K^{-1}.mol^{-1}},
\]
d'écart-type $R \times 0.103 = \qty{0.86}{J.K^{-1}.mol^{-1}}$ et d'intervalle à 95\,\% $\pm\qty{2.7}{J.K^{-1}.mol^{-1}}$.
\end{solution}

\begin{solution}{exo:b3:rate-theories:11}
À partir de $\ln k = \ln(k_BT/h) + \Delta^{\ddagger}S^\circ/R - \Delta^{\ddagger}H^\circ/RT$,
\[
  \frac{\dd\ln k}{\dd T} = \frac1T + \frac{\Delta^{\ddagger}H^\circ}{RT^2}, \qquad
  E_a = RT^2\frac{\dd\ln k}{\dd T} = \Delta^{\ddagger}H^\circ + RT .
\] Alors $A = k\eu^{E_a/RT} = (k_BT/h)\eu^{\Delta^{\ddagger}S^\circ/R}\eu^{-\Delta^{\ddagger}H^\circ/RT}\eu^{(\Delta^{\ddagger}H^\circ + RT)/RT} =
\eu\,(k_BT/h)\eu^{\Delta^{\ddagger}S^\circ/R}$. Avec $\Delta^{\ddagger}S^\circ = 0$~: $A = 2.718 \times \num{6.21e12} = \qty{1.7e13}{s^{-1}}$,
l'ordre de grandeur des facteurs observés~; des valeurs supérieures
signifient un $\Delta^{\ddagger}S^\circ$ positif, un état de transition
lâche.
\end{solution}

\begin{solution}{exo:b3:rate-theories:12}
Avec des \omterm{def:b3:partition-functions:partition-function}{fonctions de partition moléculaires} par unité de volume~:
$q^{\mathrm T}/V = (2\pi mk_BT)^{3/2}/h^3$ pour
A, B et le complexe (de masse $m_{\mathrm A} + m_{\mathrm B}$), et pour le complexe $q^{\mathrm R} = 8\pi^2Ik_BT/h^2$
($\sigma = 1$). Le rapport des facteurs de translation vaut
$h^3/(2\pi\mu k_BT)^{3/2}$, donc
\[
  k = \frac{k_BT}{h}\cdot\frac{h^3}{(2\pi\mu k_BT)^{3/2}}\cdot\frac{8\pi^2\mu d^2k_BT}{h^2}\,\eu^{-E_0/RT}
    = \frac{8\pi^2d^2(k_BT)^{1/2}}{(2\pi)^{3/2}\mu^{1/2}}\eu^{-E_0/RT}
    = \pi d^2\sqrt{\frac{8k_BT}{\pi\mu}}\,\eu^{-E_0/RT},
\]
par paire de molécules~; multiplier par $N_A$ donne le
\cref{thm:b3:rate-theories:collision}
avec $\sigma = \pi d^2$.
\end{solution}

\begin{solution}{pb:b3:rate-theories:1}
\textbf{1.} $\mu_{\mathrm{CH}} = 12 \times 1.00783/13.00783 = \qty{0.9297}{u}$~; $\mu_{\mathrm{CD}} = 12 \times 2.01410/14.01410 =
\qty{1.7246}{u}$.
\textbf{2.} $2917/\sqrt{1.7246/0.9297} = 2917/1.3620 = \qty{2142}{cm^{-1}}$,
1.5\,\% au-dessus de la valeur
mesurée, 2109 (anharmonicité et couplage des quatre liaisons de
\ce{CD4}).
\textbf{3.} $\frac12 \times 2917 = \qty{1458.5}{cm^{-1}}$ et $\frac12 \times 2109 = \qty{1054.5}{cm^{-1}}$~;
différence
\qty{404}{cm^{-1}} $= \qty{4.83}{kJ/mol}$.
\textbf{4.} Elle devient la coordonnée réactionnelle~: son \omterm{def:b3:quantum-model-systems:oscillator}{énergie de
point zéro} disparaît.
\textbf{5.} La barrière est plus haute pour D de \qty{4.83}{kJ/mol}.
\textbf{6.} $\exp(1.43878 \times 404/310.15) = \eu^{1.874} = 6.5$.
\textbf{7.} $\exp(1.43878 \times 404/298.15) = 7.0$.
\textbf{8.} Dans un état de transition réel, une partie de l'\omterm{def:b3:quantum-model-systems:oscillator}{énergie de
point zéro} de l'élongation subsiste dans d'autres modes, ce qui diminue
l'effet~; les seules \omterm{def:b3:quantum-model-systems:oscillator}{énergies de point zéro} ne peuvent pas dépasser
cette valeur (l'\omterm{def:b3:quantum-model-systems:tunnelling}{effet tunnel} le peut).
\textbf{9.} Peu~: ce sont des effets secondaires, de l'ordre de 10\,\% par
deutérium au plus.
\textbf{10.} Que la liaison C--H est rompue dans l'étape cinétiquement
déterminante, l'élongation étant entièrement convertie en coordonnée
réactionnelle.
\textbf{11.} Élimination d'ordre 1~: $t_{1/2} = \ln2/k_{\mathrm{total}}$,
inversement proportionnel à cette constante.
\textbf{12.} $k_{\mathrm{total}}(\mathrm H) = 0.75\,k_{\mathrm{total}}(\mathrm H) + 0.25\,k_{\mathrm{total}}(\mathrm H)$,
déméthylation et autres voies.
\textbf{13.} Divisée par 6.5.
\textbf{14.} $0.25 + 0.75/6.52 = 0.365$.
\textbf{15.} $5.0/0.365 = \qty{13.7}{h}$.
\textbf{16.} 2.7.
\textbf{17.} 6.5, l'effet isotopique entier.
\textbf{18.} Des doses plus faibles ou moins fréquentes, et des
concentrations sanguines plus régulières.
\textbf{19.} $\ln(k/T)$ en fonction de $1/T$.
\textbf{20.} $\Delta^{\ddagger}H^\circ = 61.9 \pm \qty{0.3}{kJ/mol}$, $\Delta^{\ddagger}S^\circ = -60.1 \pm \qty{0.9}{J.K^{-1}.mol^{-1}}$.
\textbf{21.} $\Delta^{\ddagger}H^\circ = 66.8 \pm \qty{0.2}{kJ/mol}$, $\Delta^{\ddagger}S^\circ = -60.3 \pm \qty{0.7}{J.K^{-1}.mol^{-1}}$.
\textbf{22.} $\Delta\Delta^{\ddagger}H^\circ = 4.88 \pm \qty{0.33}{kJ/mol}$ ($\sqrt{0.26^2 + 0.21^2}$).
\textbf{23.} Oui~: 4.83 est bien à moins d'un écart-type. L'effet
isotopique s'explique entièrement par les \omterm{def:b3:quantum-model-systems:oscillator}{énergies de point zéro}, sans
signe d'\omterm{def:b3:quantum-model-systems:tunnelling}{effet tunnel}.
\textbf{24.} Elles sont égales aux erreurs près~: l'état de transition a
la même structure pour les deux \omterm{def:b3:rovibrational-spectroscopy:rotational-constant}{isotopologues}, comme on l'a supposé dans
la partie I.
\textbf{25.} \textbf{$t_{1/2}(\mathrm D)/t_{1/2}(\mathrm H) \approx 2.7$
à \qty{37}{\celsius}}~: \qty{13.7}{h} contre \qty{5.0}{h}.
\end{solution}
